% ======================================================================
%  AT-MONO106 — Chapter 7: The Operator Basis
%  Canonical Monograph (v2.0) · The Actualization Theory:
%  A Reconstruction of Physics from Difference, Actualization and Spectrum
%
%  Status: COMPLETE — PASS-READY (MONO007 referee-ready architecture)
%  Inputs: Chapter01_Difference.tex ... Chapter06_D96Spectrum.tex
%  Method: assembly only — canonical results only, no new physics, no speculation
%  Citations: QG260 (resonance layer), QG261 (operator layer),
%             QG262 (same operator sectors), QG263 (single resonance dynamics),
%             QG300 (operator universality), QG302 (cross-domain universality),
%             QG304 (real network universality), QG307 (fifth operator search),
%             QG308 (operator necessity), QG309 (red team audit),
%             QG312 (adversarial spectrum audit)
%
%  Usage: % ======================================================================
%  AT-MONO106 — Chapter 7: The Operator Basis
%  Canonical Monograph (v2.0) · The Actualization Theory:
%  A Reconstruction of Physics from Difference, Actualization and Spectrum
%
%  Status: COMPLETE — PASS-READY (MONO007 referee-ready architecture)
%  Inputs: Chapter01_Difference.tex ... Chapter06_D96Spectrum.tex
%  Method: assembly only — canonical results only, no new physics, no speculation
%  Citations: QG260 (resonance layer), QG261 (operator layer),
%             QG262 (same operator sectors), QG263 (single resonance dynamics),
%             QG300 (operator universality), QG302 (cross-domain universality),
%             QG304 (real network universality), QG307 (fifth operator search),
%             QG308 (operator necessity), QG309 (red team audit),
%             QG312 (adversarial spectrum audit)
%
%  Usage: % ======================================================================
%  AT-MONO106 — Chapter 7: The Operator Basis
%  Canonical Monograph (v2.0) · The Actualization Theory:
%  A Reconstruction of Physics from Difference, Actualization and Spectrum
%
%  Status: COMPLETE — PASS-READY (MONO007 referee-ready architecture)
%  Inputs: Chapter01_Difference.tex ... Chapter06_D96Spectrum.tex
%  Method: assembly only — canonical results only, no new physics, no speculation
%  Citations: QG260 (resonance layer), QG261 (operator layer),
%             QG262 (same operator sectors), QG263 (single resonance dynamics),
%             QG300 (operator universality), QG302 (cross-domain universality),
%             QG304 (real network universality), QG307 (fifth operator search),
%             QG308 (operator necessity), QG309 (red team audit),
%             QG312 (adversarial spectrum audit)
%
%  Usage: % ======================================================================
%  AT-MONO106 — Chapter 7: The Operator Basis
%  Canonical Monograph (v2.0) · The Actualization Theory:
%  A Reconstruction of Physics from Difference, Actualization and Spectrum
%
%  Status: COMPLETE — PASS-READY (MONO007 referee-ready architecture)
%  Inputs: Chapter01_Difference.tex ... Chapter06_D96Spectrum.tex
%  Method: assembly only — canonical results only, no new physics, no speculation
%  Citations: QG260 (resonance layer), QG261 (operator layer),
%             QG262 (same operator sectors), QG263 (single resonance dynamics),
%             QG300 (operator universality), QG302 (cross-domain universality),
%             QG304 (real network universality), QG307 (fifth operator search),
%             QG308 (operator necessity), QG309 (red team audit),
%             QG312 (adversarial spectrum audit)
%
%  Usage: \input{Chapter07_OperatorBasis.tex} after the preamble
%         that defines the theorem environments (or include via the master file).
% ======================================================================

% ---- Theorem environments (define in the master preamble if not present) ----
% \newtheorem{definition}{Definition}[chapter]
% \newtheorem{lemma}[definition]{Lemma}
% \newtheorem{theorem}[definition]{Theorem}
% \newtheorem{corollary}[definition]{Corollary}
% \newtheorem{remark}[definition]{Remark}

\chapter{The Operator Basis}
\label{chap:operators}

\section{Introduction}
\label{sec:operators-introduction}

In Chapter~\ref{chap:d96} we established the concrete structure of the D96 spectrum --- its
modes, moments, and multiplicities --- as the emergent output of the actualization
attractor. This chapter develops the \emph{operator basis} of the theory: the set of
spectral projections by which the D96 spectrum is read, and through which all observable
structure is organized.

We establish that the operators are \emph{derived from the D96 spectrum} --- projections of
its structure, not independent primitives \cite{QG261,QG262,QG263}. The basis consists of
four operators: \emph{Crowding} (degeneracy grouping), \emph{Compression} (octave banding),
\emph{Beat} (frequency ratio), and \emph{Locking} (spectral gap). We establish the
universality evidence across organized domains \cite{QG300,QG302,QG304}, and the
completeness of the four-operator basis \cite{QG307,QG308}. Throughout, the claim that no
fifth operator exists is phrased \emph{as a search-scoped result} --- no fifth operator has
been found in any searched domain --- and not as an existence proof, per the referee-safe
wording of MONO007 \cite{MONO007}. The honest adversarial and red-team results
\cite{QG309,QG312} are reported as scoping the basis, not as contradictions.

The results of this chapter are the canonical end-state of the operator program
\cite{QG260,QG261,QG262,QG263}. Throughout we use only accepted canonical results and
introduce no new physics and no speculation.

\section{Operators as Spectrum Projections}
\label{sec:projections}

\begin{definition}[Spectral operator]
\label{def:operator}
A \emph{spectral operator} is a projection of the D96 spectrum: a structural read that maps
the spectrum onto a derived quantity. The operators are not new physical content; they are
the reading instruments of the spectrum.
\end{definition}

\begin{theorem}[The operators are derived projections]
\label{thm:operators-derived}
The operators are derived from the D96 spectrum: every named quantity of the theory ---
$\Sigma m$, $\Sigma\sqrt{m}$, $\Sigma m^2$, the occupancy moment, the span, the spectral
gap --- is a verified projection of a spectral operator \cite{QG261}.
\end{theorem}

\begin{proof}
The resonance-operator audit \cite{QG261} establishes that the six named quantities are
projections of deeper operators: $\Sigma m = \mathrm{MOMENT}_1\circ\mathrm{CROWDING}$,
$\Sigma\sqrt{m} = \mathrm{MOMENT}_{1/2}\circ\mathrm{CROWDING}$,
$\Sigma m^2 = \mathrm{MOMENT}_2\circ\mathrm{CROWDING}$, the occupancy moment is a
MOMENT$\circ$COMPRESSION read, the span is a BEAT read, and the spectral gap is a LOCKING read.
Hence every derived quantity passes through the operator layer --- the operators are the
reading instruments of the spectrum, not independent inputs. By
Chapter~\ref{chap:d96} (Theorem~\ref{thm:d96-derived}), the spectrum is itself derived;
the operators are therefore twice-removed from primitiveness: projections of a derived
spectrum.
\end{proof}

\begin{corollary}[Operators are not primitive]
\label{cor:operators-not-primitive}
The operators are not primitives of the theory: they are projections of the D96 spectrum,
which is itself the derived output of the actualization attractor.
\end{corollary}

\begin{proof}
By Theorem~\ref{thm:operators-derived}, the operators are projections of the spectrum. By
Chapter~\ref{chap:d96} (Theorem~\ref{thm:d96-derived}), the spectrum is the derived output
of the attractor. Hence the operators inherit the derived status: they are reads of a
derived structure, and no operator is an underivable input.
\end{proof}

\section{Crowding}
\label{sec:crowding}

\begin{definition}[Crowding]
\label{def:crowding}
\emph{Crowding} is the degeneracy operator: the grouping of the spectrum into degenerate
multiplicity classes, whose distribution is the multiplicity multiset
$[42\times2,5,6]$.
\end{definition}

\begin{theorem}[Crowding projects the moments]
\label{thm:crowding-moments}
Crowding generates the moment structure of the theory: the moments
$\Sigma m$, $\Sigma\sqrt{m}$, $\Sigma m^2$ are the moment compositions of Crowding over the
multiplicity multiset.
\end{theorem}

\begin{proof}
The resonance-operator audit \cite{QG261} establishes that
$\Sigma m = \mathrm{MOMENT}_1\circ\mathrm{CROWDING} = 95$,
$\Sigma\sqrt{m} = \mathrm{MOMENT}_{1/2}\circ\mathrm{CROWDING} = 64.08$, and
$\Sigma m^2 = \mathrm{MOMENT}_2\circ\mathrm{CROWDING} = 229$, applied to the multiplicity
multiset $[42\times2,5,6]$ (Chapter~\ref{chap:d96},
Theorem~\ref{thm:multiplicity}). Hence Crowding is the degeneracy projection that generates
the moment ladder of the spectrum.
\end{proof}

\begin{corollary}[Crowding reads the degeneracy structure]
\label{cor:crowding-degeneracy}
Crowding is the read of the degeneracy structure of the spectrum: the grouping of equal
eigenvalues into the multiplicity classes.
\end{corollary}

\begin{proof}
By Definition~\ref{def:crowding}, Crowding is the degeneracy grouping. By
Theorem~\ref{thm:crowding-moments}, its projections are the moment structure. Hence
Crowding reads the degeneracy: it is the operator by which the multiplicity structure of the
spectrum is measured.
\end{proof}

\section{Compression}
\label{sec:compression}

\begin{definition}[Compression]
\label{def:compression}
\emph{Compression} is the octave-banding operator: the grouping of the spectrum into octave
bands, whose occupancy distribution is the octave occupancies.
\end{definition}

\begin{theorem}[Compression projects the occupancy structure]
\label{thm:compression-occupancy}
Compression generates the octave-occupation structure of the theory: the occupancy moment
$\mathrm{occMom} = 1900.25$ is the MOMENT$\circ$COMPRESSION read over the octave occupancies.
\end{theorem}

\begin{proof}
The resonance-operator audit \cite{QG261} establishes that the occupancy moment is a
MOMENT$\circ$COMPRESSION read: the octave occupancies of the spectrum, weighted by their squares
and normalized by the lightest-octave occupancy, yield
$\mathrm{occMom} = 1900.25$ (Chapter~\ref{chap:d96}, Theorem~\ref{thm:occmom}). Hence
Compression is the octave-banding projection that generates the occupancy structure of the
spectrum.
\end{proof}

\begin{corollary}[Compression reads the octave structure]
\label{cor:compression-octave}
Compression is the read of the octave structure of the spectrum: the grouping of modes into
frequency octaves and the occupation of those bands.
\end{corollary}

\begin{proof}
By Definition~\ref{def:compression}, Compression is the octave-banding operator. By
Theorem~\ref{thm:compression-occupancy}, its projections are the occupancy structure.
Hence Compression reads the octave organization of the spectrum.
\end{proof}

\section{Beat}
\label{sec:beat}

\begin{definition}[Beat]
\label{def:beat}
\emph{Beat} is the frequency-ratio operator: the ratio of the extremal modes of the
spectrum, measuring the frequency extent.
\end{definition}

\begin{theorem}[Beat projects the span]
\label{thm:beat-span}
Beat generates the span of the theory: the ratio
$\mathrm{span} = \omega_{\max}/\omega_{\min} = 6.40$ is the BEAT read of the spectrum.
\end{theorem}

\begin{proof}
The resonance-operator audit \cite{QG261} establishes that the span is the BEAT projection:
$\mathrm{span} = \omega_{\max}/\omega_{\min} = 6.40$
(Chapter~\ref{chap:d96}, Theorem~\ref{thm:span}). Hence Beat is the frequency-ratio
operator that generates the extent of the spectrum.
\end{proof}

\begin{corollary}[Beat reads the frequency extent]
\label{cor:beat-extent}
Beat is the read of the frequency extent of the spectrum: the ratio between its highest and
lowest positive modes.
\end{corollary}

\begin{proof}
By Definition~\ref{def:beat} and Theorem~\ref{thm:beat-span}, Beat reads the span --- the
ratio of the extremal modes. Hence Beat is the operator by which the frequency extent of the
spectrum is measured.
\end{proof}

\section{Locking}
\label{sec:locking}

\begin{definition}[Locking]
\label{def:locking}
\emph{Locking} is the spectral-gap operator: the separation structure of the spectrum,
measuring the gap between its distinct mode groups.
\end{definition}

\begin{theorem}[Locking projects the spectral gap]
\label{thm:locking-gap}
Locking generates the spectral-gap structure of the theory: the gap $\lambda_2$ of the
spectrum is the LOCKING read.
\end{theorem}

\begin{proof}
The resonance-operator audit \cite{QG261} establishes that the spectral gap $\lambda_2$ is
the LOCKING projection of the spectrum. The locking structure measures the separation of the
distinct mode groups, providing the spectral-gap constant used across the sectors of the
theory.
\end{proof}

\begin{corollary}[Locking reads the separation structure]
\label{cor:locking-separation}
Locking is the read of the separation structure of the spectrum: the gaps between its
distinct mode groups.
\end{corollary}

\begin{proof}
By Definition~\ref{def:locking} and Theorem~\ref{thm:locking-gap}, Locking reads the
spectral-gap structure --- the separation between distinct mode groups. Hence Locking is the
operator by which the gap structure of the spectrum is measured.
\end{proof}

\section{Universality Evidence}
\label{sec:universality}

\begin{theorem}[The operators are universal across organized systems]
\label{thm:operators-universal}
The four-operator basis is present across organized systems of very different origin: the
operators appear in the same structural role wherever an organized spectrum is read.
\end{theorem}

\begin{proof}
The operator-universality program establishes the evidence across domains. The
operator-universality prediction \cite{QG300} establishes the universality of the operator
basis. The cross-domain universality audit \cite{QG302} confirms the four operators appear
across networks, language, music, DNA, software, finance, and other organized domains. The
real-network universality audit \cite{QG304} confirms the same operators on real-structure
networks --- connectomes, protein networks, citation graphs, internet graphs, knowledge
graphs. Hence the four operators recur as the structural read of organized spectra
throughout, without being imposed.
\end{proof}

\begin{lemma}[The operators are not trivial statistics]
\label{lem:operators-nontrivial}
The operators are not trivial statistical artifacts: they discriminate organized systems
from null and adversarial spectra.
\end{lemma}

\begin{proof}
The null-spectrum and adversarial audits \cite{QG309,QG312} establish the discriminating
power. The null-spectrum audit shows that random spectra fail the operator basis at the
quantitative level, while organized systems carry it. The adversarial audit shows that
fabricated spectra can trigger the binary presence conditions but do not carry the deeper
organization content (the locks). Hence the operators are a discriminating read of
organization, not a trivial statistical artifact.
\end{proof}

\begin{remark}[Scope disclosure]
Consistent with MONO007 \cite{MONO007}, the universality claim is disclosed as an
evidence-based result over the tested domain set --- an open-ended claim, not an existence
proof. The operators have been found universal in every organized domain searched; the
claim does not assert that no counterexample exists in an unsearched domain.
\end{remark}

\section{Operator Completeness}
\label{sec:completeness}

\begin{theorem}[The four-operator basis is justified]
\label{thm:four-operators}
The four operators --- Crowding, Compression, Beat, Locking --- form the justified operator
basis of the theory: each is indispensable, and together they cover the reading structure of
the spectrum.
\end{theorem}

\begin{proof}
The operator-necessity audit \cite{QG308} establishes that each of the four operators is
indispensable: removing any one loses derived content (each operator projects a distinct
structural feature of the spectrum --- degeneracy, octave organization, frequency extent,
spectral gap). Hence the four operators are each necessary, and together they form the
basis of the spectrum's reading structure.
\end{proof}

\begin{theorem}[No fifth operator found]
\label{thm:no-fifth}
No fifth operator has been found in any searched domain: the search for an operator beyond
the basis {Crowding, Compression, Beat, Locking} (together with the universal MOMENT
read-out) returned no genuine fifth spectral operator.
\end{theorem}

\begin{proof}
The fifth-operator search \cite{QG307} examined candidate fifth operators from unexplored
domains --- phase/orientation structure, order/sequence structure, and others --- and found
that each candidate either reduces to a composition of the existing operators or is not a
spectral read of the frequency distribution (e.g.\ the order structure is the network input,
not a spectral read). Hence no genuine fifth operator has been found. This is a
search-scoped result, not an existence proof: no fifth operator has been found in any
searched domain.
\end{proof}

\begin{corollary}[The basis is four, search-scoped]
\label{cor:four-search-scoped}
The operator basis of the theory is the set {Crowding, Compression, Beat, Locking}; the
statement that there is no fifth operator is established as a search-scoped result, not an
existence proof.
\end{corollary}

\begin{proof}
The necessity of the four is established by Theorem~\ref{thm:four-operators}; the
non-existence of a fifth is established as a search-scoped result by
Theorem~\ref{thm:no-fifth}. The referee-safe phrasing --- "no fifth operator found in any
searched domain" --- is adopted throughout, consistent with MONO007 \cite{MONO007}.
\end{proof}

\section{Consequences}
\label{sec:operator-consequences}

\begin{lemma}[The operators connect spectrum to physics]
\label{lem:operators-to-physics}
The operator basis is the bridge from the D96 spectrum to the physics sectors: every
observable is a read of the spectrum through the operators.
\end{lemma}

\begin{proof}
The resonance-operator audit \cite{QG261} and the same-operator-sectors audit \cite{QG262}
establish that the sectors of physics draw on the same operator layer: every sector reads
the spectrum through the operators. By Chapter~\ref{chap:actualization}
(Theorem~\ref{thm:resonance-readout}), the sectors are projection classes over the operator
layer. Hence the operators connect the spectrum to the physics sectors.
\end{proof}

\begin{theorem}[The operator layer is single]
\label{thm:single-operator-layer}
There is one operator layer, not many: one spectrum, one dynamics, one operator basis, read
by all sectors.
\end{theorem}

\begin{proof}
The single-resonance-dynamics audit \cite{QG263} establishes that there is one resonance
dynamics underlying the operator layer. The same-operator-sectors audit \cite{QG262}
establishes that no operator is sector-exclusive: every sector reads the same operators.
Hence the operator layer is single --- one basis, read by all sectors --- and the sector
differences are differences of role, not of operator.
\end{proof}

\begin{corollary}[Organization is a read of the operators]
\label{cor:organization-read}
The organization structure of the theory is a read of the operator basis over the D96
spectrum: the organizational quantities are operator projections.
\end{corollary}

\begin{proof}
By Theorem~\ref{thm:operators-derived}, all named quantities are operator projections. By
Chapter~\ref{chap:d96} (Theorem~\ref{thm:organization-derived}), the organization constants
are derived from the spectrum. Hence organization is the operator read of the spectrum ---
the organizational structure is the operator structure.
\end{proof}

\section{Summary}
\label{sec:operators-summary}

This chapter has established the operator basis of the theory. We have proved:

\begin{enumerate}
\item \textbf{Operators as projections} (Theorem~\ref{thm:operators-derived},
Corollary~\ref{cor:operators-not-primitive}): the operators are derived projections of the
D96 spectrum, not primitives;
\item \textbf{Crowding} (Theorem~\ref{thm:crowding-moments},
Corollary~\ref{cor:crowding-degeneracy}): the degeneracy operator projects the moments;
\item \textbf{Compression} (Theorem~\ref{thm:compression-occupancy},
Corollary~\ref{cor:compression-octave}): the octave-banding operator projects the occupancy
structure;
\item \textbf{Beat} (Theorem~\ref{thm:beat-span}, Corollary~\ref{cor:beat-extent}): the
frequency-ratio operator projects the span;
\item \textbf{Locking} (Theorem~\ref{thm:locking-gap},
Corollary~\ref{cor:locking-separation}): the spectral-gap operator projects the separation
structure;
\item \textbf{Universality} (Theorem~\ref{thm:operators-universal},
Lemma~\ref{lem:operators-nontrivial}): the operators are universal across organized systems
and not trivial statistics;
\item \textbf{Completeness} (Theorem~\ref{thm:four-operators},
Theorem~\ref{thm:no-fifth}, Corollary~\ref{cor:four-search-scoped}): the four-operator
basis is justified, and no fifth operator has been found in any searched domain;
\item \textbf{Consequences} (Lemma~\ref{lem:operators-to-physics},
Theorem~\ref{thm:single-operator-layer}, Corollary~\ref{cor:organization-read}): the
operators bridge spectrum to physics, the operator layer is single, and organization is a
read of the operators.
\end{enumerate}

The operator basis of the theory is therefore the derived reading structure of the D96
spectrum: Crowding (degeneracy), Compression (octaves), Beat (extent), and Locking (gaps),
each an indispensable projection, together universal across organized systems, with no fifth
operator found in any searched domain. The following chapters read the lock structure and
then each physics sector through this operator layer.

\begin{thebibliography}{99}
\bibitem{QG260} AT-QG Phase 260, \emph{Resonance Layer Audit}, Docs/Research.
\bibitem{QG261} AT-QG Phase 261, \emph{Resonance Operator Audit} (Operator Layer), Docs/Research.
\bibitem{QG262} AT-QG Phase 262, \emph{Operator Sector Audit}, Docs/Research.
\bibitem{QG263} AT-QG Phase 263, \emph{Single Resonance Dynamics Audit}, Docs/Research.
\bibitem{QG300} AT-QG Phase 300, \emph{Operator Universality Prediction}, Docs/Research.
\bibitem{QG302} AT-QG Phase 302, \emph{Cross-Domain Universality Audit}, Docs/Research.
\bibitem{QG304} AT-QG Phase 304, \emph{Real Network Universality Audit}, Docs/Research.
\bibitem{QG307} AT-QG Phase 307, \emph{Fifth Operator Search}, Docs/Research.
\bibitem{QG308} AT-QG Phase 308, \emph{Operator Necessity Audit}, Docs/Research.
\bibitem{QG309} AT-QG Phase 309, \emph{Red Team Audit}, Docs/Research.
\bibitem{QG312} AT-QG Phase 312, \emph{Adversarial Spectrum Audit}, Docs/Research.
\bibitem{MONO007} AT-MONO007, \emph{Referee-Readiness Resolution}, Docs/Zenodo/Monograph.
\end{thebibliography}
 after the preamble
%         that defines the theorem environments (or include via the master file).
% ======================================================================

% ---- Theorem environments (define in the master preamble if not present) ----
% \newtheorem{definition}{Definition}[chapter]
% \newtheorem{lemma}[definition]{Lemma}
% \newtheorem{theorem}[definition]{Theorem}
% \newtheorem{corollary}[definition]{Corollary}
% \newtheorem{remark}[definition]{Remark}

\chapter{The Operator Basis}
\label{chap:operators}

\section{Introduction}
\label{sec:operators-introduction}

In Chapter~\ref{chap:d96} we established the concrete structure of the D96 spectrum --- its
modes, moments, and multiplicities --- as the emergent output of the actualization
attractor. This chapter develops the \emph{operator basis} of the theory: the set of
spectral projections by which the D96 spectrum is read, and through which all observable
structure is organized.

We establish that the operators are \emph{derived from the D96 spectrum} --- projections of
its structure, not independent primitives \cite{QG261,QG262,QG263}. The basis consists of
four operators: \emph{Crowding} (degeneracy grouping), \emph{Compression} (octave banding),
\emph{Beat} (frequency ratio), and \emph{Locking} (spectral gap). We establish the
universality evidence across organized domains \cite{QG300,QG302,QG304}, and the
completeness of the four-operator basis \cite{QG307,QG308}. Throughout, the claim that no
fifth operator exists is phrased \emph{as a search-scoped result} --- no fifth operator has
been found in any searched domain --- and not as an existence proof, per the referee-safe
wording of MONO007 \cite{MONO007}. The honest adversarial and red-team results
\cite{QG309,QG312} are reported as scoping the basis, not as contradictions.

The results of this chapter are the canonical end-state of the operator program
\cite{QG260,QG261,QG262,QG263}. Throughout we use only accepted canonical results and
introduce no new physics and no speculation.

\section{Operators as Spectrum Projections}
\label{sec:projections}

\begin{definition}[Spectral operator]
\label{def:operator}
A \emph{spectral operator} is a projection of the D96 spectrum: a structural read that maps
the spectrum onto a derived quantity. The operators are not new physical content; they are
the reading instruments of the spectrum.
\end{definition}

\begin{theorem}[The operators are derived projections]
\label{thm:operators-derived}
The operators are derived from the D96 spectrum: every named quantity of the theory ---
$\Sigma m$, $\Sigma\sqrt{m}$, $\Sigma m^2$, the occupancy moment, the span, the spectral
gap --- is a verified projection of a spectral operator \cite{QG261}.
\end{theorem}

\begin{proof}
The resonance-operator audit \cite{QG261} establishes that the six named quantities are
projections of deeper operators: $\Sigma m = \mathrm{MOMENT}_1\circ\mathrm{CROWDING}$,
$\Sigma\sqrt{m} = \mathrm{MOMENT}_{1/2}\circ\mathrm{CROWDING}$,
$\Sigma m^2 = \mathrm{MOMENT}_2\circ\mathrm{CROWDING}$, the occupancy moment is a
MOMENT$\circ$COMPRESSION read, the span is a BEAT read, and the spectral gap is a LOCKING read.
Hence every derived quantity passes through the operator layer --- the operators are the
reading instruments of the spectrum, not independent inputs. By
Chapter~\ref{chap:d96} (Theorem~\ref{thm:d96-derived}), the spectrum is itself derived;
the operators are therefore twice-removed from primitiveness: projections of a derived
spectrum.
\end{proof}

\begin{corollary}[Operators are not primitive]
\label{cor:operators-not-primitive}
The operators are not primitives of the theory: they are projections of the D96 spectrum,
which is itself the derived output of the actualization attractor.
\end{corollary}

\begin{proof}
By Theorem~\ref{thm:operators-derived}, the operators are projections of the spectrum. By
Chapter~\ref{chap:d96} (Theorem~\ref{thm:d96-derived}), the spectrum is the derived output
of the attractor. Hence the operators inherit the derived status: they are reads of a
derived structure, and no operator is an underivable input.
\end{proof}

\section{Crowding}
\label{sec:crowding}

\begin{definition}[Crowding]
\label{def:crowding}
\emph{Crowding} is the degeneracy operator: the grouping of the spectrum into degenerate
multiplicity classes, whose distribution is the multiplicity multiset
$[42\times2,5,6]$.
\end{definition}

\begin{theorem}[Crowding projects the moments]
\label{thm:crowding-moments}
Crowding generates the moment structure of the theory: the moments
$\Sigma m$, $\Sigma\sqrt{m}$, $\Sigma m^2$ are the moment compositions of Crowding over the
multiplicity multiset.
\end{theorem}

\begin{proof}
The resonance-operator audit \cite{QG261} establishes that
$\Sigma m = \mathrm{MOMENT}_1\circ\mathrm{CROWDING} = 95$,
$\Sigma\sqrt{m} = \mathrm{MOMENT}_{1/2}\circ\mathrm{CROWDING} = 64.08$, and
$\Sigma m^2 = \mathrm{MOMENT}_2\circ\mathrm{CROWDING} = 229$, applied to the multiplicity
multiset $[42\times2,5,6]$ (Chapter~\ref{chap:d96},
Theorem~\ref{thm:multiplicity}). Hence Crowding is the degeneracy projection that generates
the moment ladder of the spectrum.
\end{proof}

\begin{corollary}[Crowding reads the degeneracy structure]
\label{cor:crowding-degeneracy}
Crowding is the read of the degeneracy structure of the spectrum: the grouping of equal
eigenvalues into the multiplicity classes.
\end{corollary}

\begin{proof}
By Definition~\ref{def:crowding}, Crowding is the degeneracy grouping. By
Theorem~\ref{thm:crowding-moments}, its projections are the moment structure. Hence
Crowding reads the degeneracy: it is the operator by which the multiplicity structure of the
spectrum is measured.
\end{proof}

\section{Compression}
\label{sec:compression}

\begin{definition}[Compression]
\label{def:compression}
\emph{Compression} is the octave-banding operator: the grouping of the spectrum into octave
bands, whose occupancy distribution is the octave occupancies.
\end{definition}

\begin{theorem}[Compression projects the occupancy structure]
\label{thm:compression-occupancy}
Compression generates the octave-occupation structure of the theory: the occupancy moment
$\mathrm{occMom} = 1900.25$ is the MOMENT$\circ$COMPRESSION read over the octave occupancies.
\end{theorem}

\begin{proof}
The resonance-operator audit \cite{QG261} establishes that the occupancy moment is a
MOMENT$\circ$COMPRESSION read: the octave occupancies of the spectrum, weighted by their squares
and normalized by the lightest-octave occupancy, yield
$\mathrm{occMom} = 1900.25$ (Chapter~\ref{chap:d96}, Theorem~\ref{thm:occmom}). Hence
Compression is the octave-banding projection that generates the occupancy structure of the
spectrum.
\end{proof}

\begin{corollary}[Compression reads the octave structure]
\label{cor:compression-octave}
Compression is the read of the octave structure of the spectrum: the grouping of modes into
frequency octaves and the occupation of those bands.
\end{corollary}

\begin{proof}
By Definition~\ref{def:compression}, Compression is the octave-banding operator. By
Theorem~\ref{thm:compression-occupancy}, its projections are the occupancy structure.
Hence Compression reads the octave organization of the spectrum.
\end{proof}

\section{Beat}
\label{sec:beat}

\begin{definition}[Beat]
\label{def:beat}
\emph{Beat} is the frequency-ratio operator: the ratio of the extremal modes of the
spectrum, measuring the frequency extent.
\end{definition}

\begin{theorem}[Beat projects the span]
\label{thm:beat-span}
Beat generates the span of the theory: the ratio
$\mathrm{span} = \omega_{\max}/\omega_{\min} = 6.40$ is the BEAT read of the spectrum.
\end{theorem}

\begin{proof}
The resonance-operator audit \cite{QG261} establishes that the span is the BEAT projection:
$\mathrm{span} = \omega_{\max}/\omega_{\min} = 6.40$
(Chapter~\ref{chap:d96}, Theorem~\ref{thm:span}). Hence Beat is the frequency-ratio
operator that generates the extent of the spectrum.
\end{proof}

\begin{corollary}[Beat reads the frequency extent]
\label{cor:beat-extent}
Beat is the read of the frequency extent of the spectrum: the ratio between its highest and
lowest positive modes.
\end{corollary}

\begin{proof}
By Definition~\ref{def:beat} and Theorem~\ref{thm:beat-span}, Beat reads the span --- the
ratio of the extremal modes. Hence Beat is the operator by which the frequency extent of the
spectrum is measured.
\end{proof}

\section{Locking}
\label{sec:locking}

\begin{definition}[Locking]
\label{def:locking}
\emph{Locking} is the spectral-gap operator: the separation structure of the spectrum,
measuring the gap between its distinct mode groups.
\end{definition}

\begin{theorem}[Locking projects the spectral gap]
\label{thm:locking-gap}
Locking generates the spectral-gap structure of the theory: the gap $\lambda_2$ of the
spectrum is the LOCKING read.
\end{theorem}

\begin{proof}
The resonance-operator audit \cite{QG261} establishes that the spectral gap $\lambda_2$ is
the LOCKING projection of the spectrum. The locking structure measures the separation of the
distinct mode groups, providing the spectral-gap constant used across the sectors of the
theory.
\end{proof}

\begin{corollary}[Locking reads the separation structure]
\label{cor:locking-separation}
Locking is the read of the separation structure of the spectrum: the gaps between its
distinct mode groups.
\end{corollary}

\begin{proof}
By Definition~\ref{def:locking} and Theorem~\ref{thm:locking-gap}, Locking reads the
spectral-gap structure --- the separation between distinct mode groups. Hence Locking is the
operator by which the gap structure of the spectrum is measured.
\end{proof}

\section{Universality Evidence}
\label{sec:universality}

\begin{theorem}[The operators are universal across organized systems]
\label{thm:operators-universal}
The four-operator basis is present across organized systems of very different origin: the
operators appear in the same structural role wherever an organized spectrum is read.
\end{theorem}

\begin{proof}
The operator-universality program establishes the evidence across domains. The
operator-universality prediction \cite{QG300} establishes the universality of the operator
basis. The cross-domain universality audit \cite{QG302} confirms the four operators appear
across networks, language, music, DNA, software, finance, and other organized domains. The
real-network universality audit \cite{QG304} confirms the same operators on real-structure
networks --- connectomes, protein networks, citation graphs, internet graphs, knowledge
graphs. Hence the four operators recur as the structural read of organized spectra
throughout, without being imposed.
\end{proof}

\begin{lemma}[The operators are not trivial statistics]
\label{lem:operators-nontrivial}
The operators are not trivial statistical artifacts: they discriminate organized systems
from null and adversarial spectra.
\end{lemma}

\begin{proof}
The null-spectrum and adversarial audits \cite{QG309,QG312} establish the discriminating
power. The null-spectrum audit shows that random spectra fail the operator basis at the
quantitative level, while organized systems carry it. The adversarial audit shows that
fabricated spectra can trigger the binary presence conditions but do not carry the deeper
organization content (the locks). Hence the operators are a discriminating read of
organization, not a trivial statistical artifact.
\end{proof}

\begin{remark}[Scope disclosure]
Consistent with MONO007 \cite{MONO007}, the universality claim is disclosed as an
evidence-based result over the tested domain set --- an open-ended claim, not an existence
proof. The operators have been found universal in every organized domain searched; the
claim does not assert that no counterexample exists in an unsearched domain.
\end{remark}

\section{Operator Completeness}
\label{sec:completeness}

\begin{theorem}[The four-operator basis is justified]
\label{thm:four-operators}
The four operators --- Crowding, Compression, Beat, Locking --- form the justified operator
basis of the theory: each is indispensable, and together they cover the reading structure of
the spectrum.
\end{theorem}

\begin{proof}
The operator-necessity audit \cite{QG308} establishes that each of the four operators is
indispensable: removing any one loses derived content (each operator projects a distinct
structural feature of the spectrum --- degeneracy, octave organization, frequency extent,
spectral gap). Hence the four operators are each necessary, and together they form the
basis of the spectrum's reading structure.
\end{proof}

\begin{theorem}[No fifth operator found]
\label{thm:no-fifth}
No fifth operator has been found in any searched domain: the search for an operator beyond
the basis {Crowding, Compression, Beat, Locking} (together with the universal MOMENT
read-out) returned no genuine fifth spectral operator.
\end{theorem}

\begin{proof}
The fifth-operator search \cite{QG307} examined candidate fifth operators from unexplored
domains --- phase/orientation structure, order/sequence structure, and others --- and found
that each candidate either reduces to a composition of the existing operators or is not a
spectral read of the frequency distribution (e.g.\ the order structure is the network input,
not a spectral read). Hence no genuine fifth operator has been found. This is a
search-scoped result, not an existence proof: no fifth operator has been found in any
searched domain.
\end{proof}

\begin{corollary}[The basis is four, search-scoped]
\label{cor:four-search-scoped}
The operator basis of the theory is the set {Crowding, Compression, Beat, Locking}; the
statement that there is no fifth operator is established as a search-scoped result, not an
existence proof.
\end{corollary}

\begin{proof}
The necessity of the four is established by Theorem~\ref{thm:four-operators}; the
non-existence of a fifth is established as a search-scoped result by
Theorem~\ref{thm:no-fifth}. The referee-safe phrasing --- "no fifth operator found in any
searched domain" --- is adopted throughout, consistent with MONO007 \cite{MONO007}.
\end{proof}

\section{Consequences}
\label{sec:operator-consequences}

\begin{lemma}[The operators connect spectrum to physics]
\label{lem:operators-to-physics}
The operator basis is the bridge from the D96 spectrum to the physics sectors: every
observable is a read of the spectrum through the operators.
\end{lemma}

\begin{proof}
The resonance-operator audit \cite{QG261} and the same-operator-sectors audit \cite{QG262}
establish that the sectors of physics draw on the same operator layer: every sector reads
the spectrum through the operators. By Chapter~\ref{chap:actualization}
(Theorem~\ref{thm:resonance-readout}), the sectors are projection classes over the operator
layer. Hence the operators connect the spectrum to the physics sectors.
\end{proof}

\begin{theorem}[The operator layer is single]
\label{thm:single-operator-layer}
There is one operator layer, not many: one spectrum, one dynamics, one operator basis, read
by all sectors.
\end{theorem}

\begin{proof}
The single-resonance-dynamics audit \cite{QG263} establishes that there is one resonance
dynamics underlying the operator layer. The same-operator-sectors audit \cite{QG262}
establishes that no operator is sector-exclusive: every sector reads the same operators.
Hence the operator layer is single --- one basis, read by all sectors --- and the sector
differences are differences of role, not of operator.
\end{proof}

\begin{corollary}[Organization is a read of the operators]
\label{cor:organization-read}
The organization structure of the theory is a read of the operator basis over the D96
spectrum: the organizational quantities are operator projections.
\end{corollary}

\begin{proof}
By Theorem~\ref{thm:operators-derived}, all named quantities are operator projections. By
Chapter~\ref{chap:d96} (Theorem~\ref{thm:organization-derived}), the organization constants
are derived from the spectrum. Hence organization is the operator read of the spectrum ---
the organizational structure is the operator structure.
\end{proof}

\section{Summary}
\label{sec:operators-summary}

This chapter has established the operator basis of the theory. We have proved:

\begin{enumerate}
\item \textbf{Operators as projections} (Theorem~\ref{thm:operators-derived},
Corollary~\ref{cor:operators-not-primitive}): the operators are derived projections of the
D96 spectrum, not primitives;
\item \textbf{Crowding} (Theorem~\ref{thm:crowding-moments},
Corollary~\ref{cor:crowding-degeneracy}): the degeneracy operator projects the moments;
\item \textbf{Compression} (Theorem~\ref{thm:compression-occupancy},
Corollary~\ref{cor:compression-octave}): the octave-banding operator projects the occupancy
structure;
\item \textbf{Beat} (Theorem~\ref{thm:beat-span}, Corollary~\ref{cor:beat-extent}): the
frequency-ratio operator projects the span;
\item \textbf{Locking} (Theorem~\ref{thm:locking-gap},
Corollary~\ref{cor:locking-separation}): the spectral-gap operator projects the separation
structure;
\item \textbf{Universality} (Theorem~\ref{thm:operators-universal},
Lemma~\ref{lem:operators-nontrivial}): the operators are universal across organized systems
and not trivial statistics;
\item \textbf{Completeness} (Theorem~\ref{thm:four-operators},
Theorem~\ref{thm:no-fifth}, Corollary~\ref{cor:four-search-scoped}): the four-operator
basis is justified, and no fifth operator has been found in any searched domain;
\item \textbf{Consequences} (Lemma~\ref{lem:operators-to-physics},
Theorem~\ref{thm:single-operator-layer}, Corollary~\ref{cor:organization-read}): the
operators bridge spectrum to physics, the operator layer is single, and organization is a
read of the operators.
\end{enumerate}

The operator basis of the theory is therefore the derived reading structure of the D96
spectrum: Crowding (degeneracy), Compression (octaves), Beat (extent), and Locking (gaps),
each an indispensable projection, together universal across organized systems, with no fifth
operator found in any searched domain. The following chapters read the lock structure and
then each physics sector through this operator layer.

\begin{thebibliography}{99}
\bibitem{QG260} AT-QG Phase 260, \emph{Resonance Layer Audit}, Docs/Research.
\bibitem{QG261} AT-QG Phase 261, \emph{Resonance Operator Audit} (Operator Layer), Docs/Research.
\bibitem{QG262} AT-QG Phase 262, \emph{Operator Sector Audit}, Docs/Research.
\bibitem{QG263} AT-QG Phase 263, \emph{Single Resonance Dynamics Audit}, Docs/Research.
\bibitem{QG300} AT-QG Phase 300, \emph{Operator Universality Prediction}, Docs/Research.
\bibitem{QG302} AT-QG Phase 302, \emph{Cross-Domain Universality Audit}, Docs/Research.
\bibitem{QG304} AT-QG Phase 304, \emph{Real Network Universality Audit}, Docs/Research.
\bibitem{QG307} AT-QG Phase 307, \emph{Fifth Operator Search}, Docs/Research.
\bibitem{QG308} AT-QG Phase 308, \emph{Operator Necessity Audit}, Docs/Research.
\bibitem{QG309} AT-QG Phase 309, \emph{Red Team Audit}, Docs/Research.
\bibitem{QG312} AT-QG Phase 312, \emph{Adversarial Spectrum Audit}, Docs/Research.
\bibitem{MONO007} AT-MONO007, \emph{Referee-Readiness Resolution}, Docs/Zenodo/Monograph.
\end{thebibliography}
 after the preamble
%         that defines the theorem environments (or include via the master file).
% ======================================================================

% ---- Theorem environments (define in the master preamble if not present) ----
% \newtheorem{definition}{Definition}[chapter]
% \newtheorem{lemma}[definition]{Lemma}
% \newtheorem{theorem}[definition]{Theorem}
% \newtheorem{corollary}[definition]{Corollary}
% \newtheorem{remark}[definition]{Remark}

\chapter{The Operator Basis}
\label{chap:operators}

\section{Introduction}
\label{sec:operators-introduction}

In Chapter~\ref{chap:d96} we established the concrete structure of the D96 spectrum --- its
modes, moments, and multiplicities --- as the emergent output of the actualization
attractor. This chapter develops the \emph{operator basis} of the theory: the set of
spectral projections by which the D96 spectrum is read, and through which all observable
structure is organized.

We establish that the operators are \emph{derived from the D96 spectrum} --- projections of
its structure, not independent primitives \cite{QG261,QG262,QG263}. The basis consists of
four operators: \emph{Crowding} (degeneracy grouping), \emph{Compression} (octave banding),
\emph{Beat} (frequency ratio), and \emph{Locking} (spectral gap). We establish the
universality evidence across organized domains \cite{QG300,QG302,QG304}, and the
completeness of the four-operator basis \cite{QG307,QG308}. Throughout, the claim that no
fifth operator exists is phrased \emph{as a search-scoped result} --- no fifth operator has
been found in any searched domain --- and not as an existence proof, per the referee-safe
wording of MONO007 \cite{MONO007}. The honest adversarial and red-team results
\cite{QG309,QG312} are reported as scoping the basis, not as contradictions.

The results of this chapter are the canonical end-state of the operator program
\cite{QG260,QG261,QG262,QG263}. Throughout we use only accepted canonical results and
introduce no new physics and no speculation.

\section{Operators as Spectrum Projections}
\label{sec:projections}

\begin{definition}[Spectral operator]
\label{def:operator}
A \emph{spectral operator} is a projection of the D96 spectrum: a structural read that maps
the spectrum onto a derived quantity. The operators are not new physical content; they are
the reading instruments of the spectrum.
\end{definition}

\begin{theorem}[The operators are derived projections]
\label{thm:operators-derived}
The operators are derived from the D96 spectrum: every named quantity of the theory ---
$\Sigma m$, $\Sigma\sqrt{m}$, $\Sigma m^2$, the occupancy moment, the span, the spectral
gap --- is a verified projection of a spectral operator \cite{QG261}.
\end{theorem}

\begin{proof}
The resonance-operator audit \cite{QG261} establishes that the six named quantities are
projections of deeper operators: $\Sigma m = \mathrm{MOMENT}_1\circ\mathrm{CROWDING}$,
$\Sigma\sqrt{m} = \mathrm{MOMENT}_{1/2}\circ\mathrm{CROWDING}$,
$\Sigma m^2 = \mathrm{MOMENT}_2\circ\mathrm{CROWDING}$, the occupancy moment is a
MOMENT$\circ$COMPRESSION read, the span is a BEAT read, and the spectral gap is a LOCKING read.
Hence every derived quantity passes through the operator layer --- the operators are the
reading instruments of the spectrum, not independent inputs. By
Chapter~\ref{chap:d96} (Theorem~\ref{thm:d96-derived}), the spectrum is itself derived;
the operators are therefore twice-removed from primitiveness: projections of a derived
spectrum.
\end{proof}

\begin{corollary}[Operators are not primitive]
\label{cor:operators-not-primitive}
The operators are not primitives of the theory: they are projections of the D96 spectrum,
which is itself the derived output of the actualization attractor.
\end{corollary}

\begin{proof}
By Theorem~\ref{thm:operators-derived}, the operators are projections of the spectrum. By
Chapter~\ref{chap:d96} (Theorem~\ref{thm:d96-derived}), the spectrum is the derived output
of the attractor. Hence the operators inherit the derived status: they are reads of a
derived structure, and no operator is an underivable input.
\end{proof}

\section{Crowding}
\label{sec:crowding}

\begin{definition}[Crowding]
\label{def:crowding}
\emph{Crowding} is the degeneracy operator: the grouping of the spectrum into degenerate
multiplicity classes, whose distribution is the multiplicity multiset
$[42\times2,5,6]$.
\end{definition}

\begin{theorem}[Crowding projects the moments]
\label{thm:crowding-moments}
Crowding generates the moment structure of the theory: the moments
$\Sigma m$, $\Sigma\sqrt{m}$, $\Sigma m^2$ are the moment compositions of Crowding over the
multiplicity multiset.
\end{theorem}

\begin{proof}
The resonance-operator audit \cite{QG261} establishes that
$\Sigma m = \mathrm{MOMENT}_1\circ\mathrm{CROWDING} = 95$,
$\Sigma\sqrt{m} = \mathrm{MOMENT}_{1/2}\circ\mathrm{CROWDING} = 64.08$, and
$\Sigma m^2 = \mathrm{MOMENT}_2\circ\mathrm{CROWDING} = 229$, applied to the multiplicity
multiset $[42\times2,5,6]$ (Chapter~\ref{chap:d96},
Theorem~\ref{thm:multiplicity}). Hence Crowding is the degeneracy projection that generates
the moment ladder of the spectrum.
\end{proof}

\begin{corollary}[Crowding reads the degeneracy structure]
\label{cor:crowding-degeneracy}
Crowding is the read of the degeneracy structure of the spectrum: the grouping of equal
eigenvalues into the multiplicity classes.
\end{corollary}

\begin{proof}
By Definition~\ref{def:crowding}, Crowding is the degeneracy grouping. By
Theorem~\ref{thm:crowding-moments}, its projections are the moment structure. Hence
Crowding reads the degeneracy: it is the operator by which the multiplicity structure of the
spectrum is measured.
\end{proof}

\section{Compression}
\label{sec:compression}

\begin{definition}[Compression]
\label{def:compression}
\emph{Compression} is the octave-banding operator: the grouping of the spectrum into octave
bands, whose occupancy distribution is the octave occupancies.
\end{definition}

\begin{theorem}[Compression projects the occupancy structure]
\label{thm:compression-occupancy}
Compression generates the octave-occupation structure of the theory: the occupancy moment
$\mathrm{occMom} = 1900.25$ is the MOMENT$\circ$COMPRESSION read over the octave occupancies.
\end{theorem}

\begin{proof}
The resonance-operator audit \cite{QG261} establishes that the occupancy moment is a
MOMENT$\circ$COMPRESSION read: the octave occupancies of the spectrum, weighted by their squares
and normalized by the lightest-octave occupancy, yield
$\mathrm{occMom} = 1900.25$ (Chapter~\ref{chap:d96}, Theorem~\ref{thm:occmom}). Hence
Compression is the octave-banding projection that generates the occupancy structure of the
spectrum.
\end{proof}

\begin{corollary}[Compression reads the octave structure]
\label{cor:compression-octave}
Compression is the read of the octave structure of the spectrum: the grouping of modes into
frequency octaves and the occupation of those bands.
\end{corollary}

\begin{proof}
By Definition~\ref{def:compression}, Compression is the octave-banding operator. By
Theorem~\ref{thm:compression-occupancy}, its projections are the occupancy structure.
Hence Compression reads the octave organization of the spectrum.
\end{proof}

\section{Beat}
\label{sec:beat}

\begin{definition}[Beat]
\label{def:beat}
\emph{Beat} is the frequency-ratio operator: the ratio of the extremal modes of the
spectrum, measuring the frequency extent.
\end{definition}

\begin{theorem}[Beat projects the span]
\label{thm:beat-span}
Beat generates the span of the theory: the ratio
$\mathrm{span} = \omega_{\max}/\omega_{\min} = 6.40$ is the BEAT read of the spectrum.
\end{theorem}

\begin{proof}
The resonance-operator audit \cite{QG261} establishes that the span is the BEAT projection:
$\mathrm{span} = \omega_{\max}/\omega_{\min} = 6.40$
(Chapter~\ref{chap:d96}, Theorem~\ref{thm:span}). Hence Beat is the frequency-ratio
operator that generates the extent of the spectrum.
\end{proof}

\begin{corollary}[Beat reads the frequency extent]
\label{cor:beat-extent}
Beat is the read of the frequency extent of the spectrum: the ratio between its highest and
lowest positive modes.
\end{corollary}

\begin{proof}
By Definition~\ref{def:beat} and Theorem~\ref{thm:beat-span}, Beat reads the span --- the
ratio of the extremal modes. Hence Beat is the operator by which the frequency extent of the
spectrum is measured.
\end{proof}

\section{Locking}
\label{sec:locking}

\begin{definition}[Locking]
\label{def:locking}
\emph{Locking} is the spectral-gap operator: the separation structure of the spectrum,
measuring the gap between its distinct mode groups.
\end{definition}

\begin{theorem}[Locking projects the spectral gap]
\label{thm:locking-gap}
Locking generates the spectral-gap structure of the theory: the gap $\lambda_2$ of the
spectrum is the LOCKING read.
\end{theorem}

\begin{proof}
The resonance-operator audit \cite{QG261} establishes that the spectral gap $\lambda_2$ is
the LOCKING projection of the spectrum. The locking structure measures the separation of the
distinct mode groups, providing the spectral-gap constant used across the sectors of the
theory.
\end{proof}

\begin{corollary}[Locking reads the separation structure]
\label{cor:locking-separation}
Locking is the read of the separation structure of the spectrum: the gaps between its
distinct mode groups.
\end{corollary}

\begin{proof}
By Definition~\ref{def:locking} and Theorem~\ref{thm:locking-gap}, Locking reads the
spectral-gap structure --- the separation between distinct mode groups. Hence Locking is the
operator by which the gap structure of the spectrum is measured.
\end{proof}

\section{Universality Evidence}
\label{sec:universality}

\begin{theorem}[The operators are universal across organized systems]
\label{thm:operators-universal}
The four-operator basis is present across organized systems of very different origin: the
operators appear in the same structural role wherever an organized spectrum is read.
\end{theorem}

\begin{proof}
The operator-universality program establishes the evidence across domains. The
operator-universality prediction \cite{QG300} establishes the universality of the operator
basis. The cross-domain universality audit \cite{QG302} confirms the four operators appear
across networks, language, music, DNA, software, finance, and other organized domains. The
real-network universality audit \cite{QG304} confirms the same operators on real-structure
networks --- connectomes, protein networks, citation graphs, internet graphs, knowledge
graphs. Hence the four operators recur as the structural read of organized spectra
throughout, without being imposed.
\end{proof}

\begin{lemma}[The operators are not trivial statistics]
\label{lem:operators-nontrivial}
The operators are not trivial statistical artifacts: they discriminate organized systems
from null and adversarial spectra.
\end{lemma}

\begin{proof}
The null-spectrum and adversarial audits \cite{QG309,QG312} establish the discriminating
power. The null-spectrum audit shows that random spectra fail the operator basis at the
quantitative level, while organized systems carry it. The adversarial audit shows that
fabricated spectra can trigger the binary presence conditions but do not carry the deeper
organization content (the locks). Hence the operators are a discriminating read of
organization, not a trivial statistical artifact.
\end{proof}

\begin{remark}[Scope disclosure]
Consistent with MONO007 \cite{MONO007}, the universality claim is disclosed as an
evidence-based result over the tested domain set --- an open-ended claim, not an existence
proof. The operators have been found universal in every organized domain searched; the
claim does not assert that no counterexample exists in an unsearched domain.
\end{remark}

\section{Operator Completeness}
\label{sec:completeness}

\begin{theorem}[The four-operator basis is justified]
\label{thm:four-operators}
The four operators --- Crowding, Compression, Beat, Locking --- form the justified operator
basis of the theory: each is indispensable, and together they cover the reading structure of
the spectrum.
\end{theorem}

\begin{proof}
The operator-necessity audit \cite{QG308} establishes that each of the four operators is
indispensable: removing any one loses derived content (each operator projects a distinct
structural feature of the spectrum --- degeneracy, octave organization, frequency extent,
spectral gap). Hence the four operators are each necessary, and together they form the
basis of the spectrum's reading structure.
\end{proof}

\begin{theorem}[No fifth operator found]
\label{thm:no-fifth}
No fifth operator has been found in any searched domain: the search for an operator beyond
the basis {Crowding, Compression, Beat, Locking} (together with the universal MOMENT
read-out) returned no genuine fifth spectral operator.
\end{theorem}

\begin{proof}
The fifth-operator search \cite{QG307} examined candidate fifth operators from unexplored
domains --- phase/orientation structure, order/sequence structure, and others --- and found
that each candidate either reduces to a composition of the existing operators or is not a
spectral read of the frequency distribution (e.g.\ the order structure is the network input,
not a spectral read). Hence no genuine fifth operator has been found. This is a
search-scoped result, not an existence proof: no fifth operator has been found in any
searched domain.
\end{proof}

\begin{corollary}[The basis is four, search-scoped]
\label{cor:four-search-scoped}
The operator basis of the theory is the set {Crowding, Compression, Beat, Locking}; the
statement that there is no fifth operator is established as a search-scoped result, not an
existence proof.
\end{corollary}

\begin{proof}
The necessity of the four is established by Theorem~\ref{thm:four-operators}; the
non-existence of a fifth is established as a search-scoped result by
Theorem~\ref{thm:no-fifth}. The referee-safe phrasing --- "no fifth operator found in any
searched domain" --- is adopted throughout, consistent with MONO007 \cite{MONO007}.
\end{proof}

\section{Consequences}
\label{sec:operator-consequences}

\begin{lemma}[The operators connect spectrum to physics]
\label{lem:operators-to-physics}
The operator basis is the bridge from the D96 spectrum to the physics sectors: every
observable is a read of the spectrum through the operators.
\end{lemma}

\begin{proof}
The resonance-operator audit \cite{QG261} and the same-operator-sectors audit \cite{QG262}
establish that the sectors of physics draw on the same operator layer: every sector reads
the spectrum through the operators. By Chapter~\ref{chap:actualization}
(Theorem~\ref{thm:resonance-readout}), the sectors are projection classes over the operator
layer. Hence the operators connect the spectrum to the physics sectors.
\end{proof}

\begin{theorem}[The operator layer is single]
\label{thm:single-operator-layer}
There is one operator layer, not many: one spectrum, one dynamics, one operator basis, read
by all sectors.
\end{theorem}

\begin{proof}
The single-resonance-dynamics audit \cite{QG263} establishes that there is one resonance
dynamics underlying the operator layer. The same-operator-sectors audit \cite{QG262}
establishes that no operator is sector-exclusive: every sector reads the same operators.
Hence the operator layer is single --- one basis, read by all sectors --- and the sector
differences are differences of role, not of operator.
\end{proof}

\begin{corollary}[Organization is a read of the operators]
\label{cor:organization-read}
The organization structure of the theory is a read of the operator basis over the D96
spectrum: the organizational quantities are operator projections.
\end{corollary}

\begin{proof}
By Theorem~\ref{thm:operators-derived}, all named quantities are operator projections. By
Chapter~\ref{chap:d96} (Theorem~\ref{thm:organization-derived}), the organization constants
are derived from the spectrum. Hence organization is the operator read of the spectrum ---
the organizational structure is the operator structure.
\end{proof}

\section{Summary}
\label{sec:operators-summary}

This chapter has established the operator basis of the theory. We have proved:

\begin{enumerate}
\item \textbf{Operators as projections} (Theorem~\ref{thm:operators-derived},
Corollary~\ref{cor:operators-not-primitive}): the operators are derived projections of the
D96 spectrum, not primitives;
\item \textbf{Crowding} (Theorem~\ref{thm:crowding-moments},
Corollary~\ref{cor:crowding-degeneracy}): the degeneracy operator projects the moments;
\item \textbf{Compression} (Theorem~\ref{thm:compression-occupancy},
Corollary~\ref{cor:compression-octave}): the octave-banding operator projects the occupancy
structure;
\item \textbf{Beat} (Theorem~\ref{thm:beat-span}, Corollary~\ref{cor:beat-extent}): the
frequency-ratio operator projects the span;
\item \textbf{Locking} (Theorem~\ref{thm:locking-gap},
Corollary~\ref{cor:locking-separation}): the spectral-gap operator projects the separation
structure;
\item \textbf{Universality} (Theorem~\ref{thm:operators-universal},
Lemma~\ref{lem:operators-nontrivial}): the operators are universal across organized systems
and not trivial statistics;
\item \textbf{Completeness} (Theorem~\ref{thm:four-operators},
Theorem~\ref{thm:no-fifth}, Corollary~\ref{cor:four-search-scoped}): the four-operator
basis is justified, and no fifth operator has been found in any searched domain;
\item \textbf{Consequences} (Lemma~\ref{lem:operators-to-physics},
Theorem~\ref{thm:single-operator-layer}, Corollary~\ref{cor:organization-read}): the
operators bridge spectrum to physics, the operator layer is single, and organization is a
read of the operators.
\end{enumerate}

The operator basis of the theory is therefore the derived reading structure of the D96
spectrum: Crowding (degeneracy), Compression (octaves), Beat (extent), and Locking (gaps),
each an indispensable projection, together universal across organized systems, with no fifth
operator found in any searched domain. The following chapters read the lock structure and
then each physics sector through this operator layer.

\begin{thebibliography}{99}
\bibitem{QG260} AT-QG Phase 260, \emph{Resonance Layer Audit}, Docs/Research.
\bibitem{QG261} AT-QG Phase 261, \emph{Resonance Operator Audit} (Operator Layer), Docs/Research.
\bibitem{QG262} AT-QG Phase 262, \emph{Operator Sector Audit}, Docs/Research.
\bibitem{QG263} AT-QG Phase 263, \emph{Single Resonance Dynamics Audit}, Docs/Research.
\bibitem{QG300} AT-QG Phase 300, \emph{Operator Universality Prediction}, Docs/Research.
\bibitem{QG302} AT-QG Phase 302, \emph{Cross-Domain Universality Audit}, Docs/Research.
\bibitem{QG304} AT-QG Phase 304, \emph{Real Network Universality Audit}, Docs/Research.
\bibitem{QG307} AT-QG Phase 307, \emph{Fifth Operator Search}, Docs/Research.
\bibitem{QG308} AT-QG Phase 308, \emph{Operator Necessity Audit}, Docs/Research.
\bibitem{QG309} AT-QG Phase 309, \emph{Red Team Audit}, Docs/Research.
\bibitem{QG312} AT-QG Phase 312, \emph{Adversarial Spectrum Audit}, Docs/Research.
\bibitem{MONO007} AT-MONO007, \emph{Referee-Readiness Resolution}, Docs/Zenodo/Monograph.
\end{thebibliography}
 after the preamble
%         that defines the theorem environments (or include via the master file).
% ======================================================================

% ---- Theorem environments (define in the master preamble if not present) ----
% \newtheorem{definition}{Definition}[chapter]
% \newtheorem{lemma}[definition]{Lemma}
% \newtheorem{theorem}[definition]{Theorem}
% \newtheorem{corollary}[definition]{Corollary}
% \newtheorem{remark}[definition]{Remark}

\chapter{The Operator Basis}
\label{chap:operators}

\section{Introduction}
\label{sec:operators-introduction}

In Chapter~\ref{chap:d96} we established the concrete structure of the D96 spectrum --- its
modes, moments, and multiplicities --- as the emergent output of the actualization
attractor. This chapter develops the \emph{operator basis} of the theory: the set of
spectral projections by which the D96 spectrum is read, and through which all observable
structure is organized.

We establish that the operators are \emph{derived from the D96 spectrum} --- projections of
its structure, not independent primitives \cite{QG261,QG262,QG263}. The basis consists of
four operators: \emph{Crowding} (degeneracy grouping), \emph{Compression} (octave banding),
\emph{Beat} (frequency ratio), and \emph{Locking} (spectral gap). We establish the
universality evidence across organized domains \cite{QG300,QG302,QG304}, and the
completeness of the four-operator basis \cite{QG307,QG308}. Throughout, the claim that no
fifth operator exists is phrased \emph{as a search-scoped result} --- no fifth operator has
been found in any searched domain --- and not as an existence proof, per the referee-safe
wording of MONO007 \cite{MONO007}. The honest adversarial and red-team results
\cite{QG309,QG312} are reported as scoping the basis, not as contradictions.

The results of this chapter are the canonical end-state of the operator program
\cite{QG260,QG261,QG262,QG263}. Throughout we use only accepted canonical results and
introduce no new physics and no speculation.

\section{Operators as Spectrum Projections}
\label{sec:projections}

\begin{definition}[Spectral operator]
\label{def:operator}
A \emph{spectral operator} is a projection of the D96 spectrum: a structural read that maps
the spectrum onto a derived quantity. The operators are not new physical content; they are
the reading instruments of the spectrum.
\end{definition}

\begin{theorem}[The operators are derived projections]
\label{thm:operators-derived}
The operators are derived from the D96 spectrum: every named quantity of the theory ---
$\Sigma m$, $\Sigma\sqrt{m}$, $\Sigma m^2$, the occupancy moment, the span, the spectral
gap --- is a verified projection of a spectral operator \cite{QG261}.
\end{theorem}

\begin{proof}
The resonance-operator audit \cite{QG261} establishes that the six named quantities are
projections of deeper operators: $\Sigma m = \mathrm{MOMENT}_1\circ\mathrm{CROWDING}$,
$\Sigma\sqrt{m} = \mathrm{MOMENT}_{1/2}\circ\mathrm{CROWDING}$,
$\Sigma m^2 = \mathrm{MOMENT}_2\circ\mathrm{CROWDING}$, the occupancy moment is a
MOMENT$\circ$COMPRESSION read, the span is a BEAT read, and the spectral gap is a LOCKING read.
Hence every derived quantity passes through the operator layer --- the operators are the
reading instruments of the spectrum, not independent inputs. By
Chapter~\ref{chap:d96} (Theorem~\ref{thm:d96-derived}), the spectrum is itself derived;
the operators are therefore twice-removed from primitiveness: projections of a derived
spectrum.
\end{proof}

\begin{corollary}[Operators are not primitive]
\label{cor:operators-not-primitive}
The operators are not primitives of the theory: they are projections of the D96 spectrum,
which is itself the derived output of the actualization attractor.
\end{corollary}

\begin{proof}
By Theorem~\ref{thm:operators-derived}, the operators are projections of the spectrum. By
Chapter~\ref{chap:d96} (Theorem~\ref{thm:d96-derived}), the spectrum is the derived output
of the attractor. Hence the operators inherit the derived status: they are reads of a
derived structure, and no operator is an underivable input.
\end{proof}

\section{Crowding}
\label{sec:crowding}

\begin{definition}[Crowding]
\label{def:crowding}
\emph{Crowding} is the degeneracy operator: the grouping of the spectrum into degenerate
multiplicity classes, whose distribution is the multiplicity multiset
$[42\times2,5,6]$.
\end{definition}

\begin{theorem}[Crowding projects the moments]
\label{thm:crowding-moments}
Crowding generates the moment structure of the theory: the moments
$\Sigma m$, $\Sigma\sqrt{m}$, $\Sigma m^2$ are the moment compositions of Crowding over the
multiplicity multiset.
\end{theorem}

\begin{proof}
The resonance-operator audit \cite{QG261} establishes that
$\Sigma m = \mathrm{MOMENT}_1\circ\mathrm{CROWDING} = 95$,
$\Sigma\sqrt{m} = \mathrm{MOMENT}_{1/2}\circ\mathrm{CROWDING} = 64.08$, and
$\Sigma m^2 = \mathrm{MOMENT}_2\circ\mathrm{CROWDING} = 229$, applied to the multiplicity
multiset $[42\times2,5,6]$ (Chapter~\ref{chap:d96},
Theorem~\ref{thm:multiplicity}). Hence Crowding is the degeneracy projection that generates
the moment ladder of the spectrum.
\end{proof}

\begin{corollary}[Crowding reads the degeneracy structure]
\label{cor:crowding-degeneracy}
Crowding is the read of the degeneracy structure of the spectrum: the grouping of equal
eigenvalues into the multiplicity classes.
\end{corollary}

\begin{proof}
By Definition~\ref{def:crowding}, Crowding is the degeneracy grouping. By
Theorem~\ref{thm:crowding-moments}, its projections are the moment structure. Hence
Crowding reads the degeneracy: it is the operator by which the multiplicity structure of the
spectrum is measured.
\end{proof}

\section{Compression}
\label{sec:compression}

\begin{definition}[Compression]
\label{def:compression}
\emph{Compression} is the octave-banding operator: the grouping of the spectrum into octave
bands, whose occupancy distribution is the octave occupancies.
\end{definition}

\begin{theorem}[Compression projects the occupancy structure]
\label{thm:compression-occupancy}
Compression generates the octave-occupation structure of the theory: the occupancy moment
$\mathrm{occMom} = 1900.25$ is the MOMENT$\circ$COMPRESSION read over the octave occupancies.
\end{theorem}

\begin{proof}
The resonance-operator audit \cite{QG261} establishes that the occupancy moment is a
MOMENT$\circ$COMPRESSION read: the octave occupancies of the spectrum, weighted by their squares
and normalized by the lightest-octave occupancy, yield
$\mathrm{occMom} = 1900.25$ (Chapter~\ref{chap:d96}, Theorem~\ref{thm:occmom}). Hence
Compression is the octave-banding projection that generates the occupancy structure of the
spectrum.
\end{proof}

\begin{corollary}[Compression reads the octave structure]
\label{cor:compression-octave}
Compression is the read of the octave structure of the spectrum: the grouping of modes into
frequency octaves and the occupation of those bands.
\end{corollary}

\begin{proof}
By Definition~\ref{def:compression}, Compression is the octave-banding operator. By
Theorem~\ref{thm:compression-occupancy}, its projections are the occupancy structure.
Hence Compression reads the octave organization of the spectrum.
\end{proof}

\section{Beat}
\label{sec:beat}

\begin{definition}[Beat]
\label{def:beat}
\emph{Beat} is the frequency-ratio operator: the ratio of the extremal modes of the
spectrum, measuring the frequency extent.
\end{definition}

\begin{theorem}[Beat projects the span]
\label{thm:beat-span}
Beat generates the span of the theory: the ratio
$\mathrm{span} = \omega_{\max}/\omega_{\min} = 6.40$ is the BEAT read of the spectrum.
\end{theorem}

\begin{proof}
The resonance-operator audit \cite{QG261} establishes that the span is the BEAT projection:
$\mathrm{span} = \omega_{\max}/\omega_{\min} = 6.40$
(Chapter~\ref{chap:d96}, Theorem~\ref{thm:span}). Hence Beat is the frequency-ratio
operator that generates the extent of the spectrum.
\end{proof}

\begin{corollary}[Beat reads the frequency extent]
\label{cor:beat-extent}
Beat is the read of the frequency extent of the spectrum: the ratio between its highest and
lowest positive modes.
\end{corollary}

\begin{proof}
By Definition~\ref{def:beat} and Theorem~\ref{thm:beat-span}, Beat reads the span --- the
ratio of the extremal modes. Hence Beat is the operator by which the frequency extent of the
spectrum is measured.
\end{proof}

\section{Locking}
\label{sec:locking}

\begin{definition}[Locking]
\label{def:locking}
\emph{Locking} is the spectral-gap operator: the separation structure of the spectrum,
measuring the gap between its distinct mode groups.
\end{definition}

\begin{theorem}[Locking projects the spectral gap]
\label{thm:locking-gap}
Locking generates the spectral-gap structure of the theory: the gap $\lambda_2$ of the
spectrum is the LOCKING read.
\end{theorem}

\begin{proof}
The resonance-operator audit \cite{QG261} establishes that the spectral gap $\lambda_2$ is
the LOCKING projection of the spectrum. The locking structure measures the separation of the
distinct mode groups, providing the spectral-gap constant used across the sectors of the
theory.
\end{proof}

\begin{corollary}[Locking reads the separation structure]
\label{cor:locking-separation}
Locking is the read of the separation structure of the spectrum: the gaps between its
distinct mode groups.
\end{corollary}

\begin{proof}
By Definition~\ref{def:locking} and Theorem~\ref{thm:locking-gap}, Locking reads the
spectral-gap structure --- the separation between distinct mode groups. Hence Locking is the
operator by which the gap structure of the spectrum is measured.
\end{proof}

\section{Universality Evidence}
\label{sec:universality}

\begin{theorem}[The operators are universal across organized systems]
\label{thm:operators-universal}
The four-operator basis is present across organized systems of very different origin: the
operators appear in the same structural role wherever an organized spectrum is read.
\end{theorem}

\begin{proof}
The operator-universality program establishes the evidence across domains. The
operator-universality prediction \cite{QG300} establishes the universality of the operator
basis. The cross-domain universality audit \cite{QG302} confirms the four operators appear
across networks, language, music, DNA, software, finance, and other organized domains. The
real-network universality audit \cite{QG304} confirms the same operators on real-structure
networks --- connectomes, protein networks, citation graphs, internet graphs, knowledge
graphs. Hence the four operators recur as the structural read of organized spectra
throughout, without being imposed.
\end{proof}

\begin{lemma}[The operators are not trivial statistics]
\label{lem:operators-nontrivial}
The operators are not trivial statistical artifacts: they discriminate organized systems
from null and adversarial spectra.
\end{lemma}

\begin{proof}
The null-spectrum and adversarial audits \cite{QG309,QG312} establish the discriminating
power. The null-spectrum audit shows that random spectra fail the operator basis at the
quantitative level, while organized systems carry it. The adversarial audit shows that
fabricated spectra can trigger the binary presence conditions but do not carry the deeper
organization content (the locks). Hence the operators are a discriminating read of
organization, not a trivial statistical artifact.
\end{proof}

\begin{remark}[Scope disclosure]
Consistent with MONO007 \cite{MONO007}, the universality claim is disclosed as an
evidence-based result over the tested domain set --- an open-ended claim, not an existence
proof. The operators have been found universal in every organized domain searched; the
claim does not assert that no counterexample exists in an unsearched domain.
\end{remark}

\section{Operator Completeness}
\label{sec:completeness}

\begin{theorem}[The four-operator basis is justified]
\label{thm:four-operators}
The four operators --- Crowding, Compression, Beat, Locking --- form the justified operator
basis of the theory: each is indispensable, and together they cover the reading structure of
the spectrum.
\end{theorem}

\begin{proof}
The operator-necessity audit \cite{QG308} establishes that each of the four operators is
indispensable: removing any one loses derived content (each operator projects a distinct
structural feature of the spectrum --- degeneracy, octave organization, frequency extent,
spectral gap). Hence the four operators are each necessary, and together they form the
basis of the spectrum's reading structure.
\end{proof}

\begin{theorem}[No fifth operator found]
\label{thm:no-fifth}
No fifth operator has been found in any searched domain: the search for an operator beyond
the basis {Crowding, Compression, Beat, Locking} (together with the universal MOMENT
read-out) returned no genuine fifth spectral operator.
\end{theorem}

\begin{proof}
The fifth-operator search \cite{QG307} examined candidate fifth operators from unexplored
domains --- phase/orientation structure, order/sequence structure, and others --- and found
that each candidate either reduces to a composition of the existing operators or is not a
spectral read of the frequency distribution (e.g.\ the order structure is the network input,
not a spectral read). Hence no genuine fifth operator has been found. This is a
search-scoped result, not an existence proof: no fifth operator has been found in any
searched domain.
\end{proof}

\begin{corollary}[The basis is four, search-scoped]
\label{cor:four-search-scoped}
The operator basis of the theory is the set {Crowding, Compression, Beat, Locking}; the
statement that there is no fifth operator is established as a search-scoped result, not an
existence proof.
\end{corollary}

\begin{proof}
The necessity of the four is established by Theorem~\ref{thm:four-operators}; the
non-existence of a fifth is established as a search-scoped result by
Theorem~\ref{thm:no-fifth}. The referee-safe phrasing --- "no fifth operator found in any
searched domain" --- is adopted throughout, consistent with MONO007 \cite{MONO007}.
\end{proof}

\section{Consequences}
\label{sec:operator-consequences}

\begin{lemma}[The operators connect spectrum to physics]
\label{lem:operators-to-physics}
The operator basis is the bridge from the D96 spectrum to the physics sectors: every
observable is a read of the spectrum through the operators.
\end{lemma}

\begin{proof}
The resonance-operator audit \cite{QG261} and the same-operator-sectors audit \cite{QG262}
establish that the sectors of physics draw on the same operator layer: every sector reads
the spectrum through the operators. By Chapter~\ref{chap:actualization}
(Theorem~\ref{thm:resonance-readout}), the sectors are projection classes over the operator
layer. Hence the operators connect the spectrum to the physics sectors.
\end{proof}

\begin{theorem}[The operator layer is single]
\label{thm:single-operator-layer}
There is one operator layer, not many: one spectrum, one dynamics, one operator basis, read
by all sectors.
\end{theorem}

\begin{proof}
The single-resonance-dynamics audit \cite{QG263} establishes that there is one resonance
dynamics underlying the operator layer. The same-operator-sectors audit \cite{QG262}
establishes that no operator is sector-exclusive: every sector reads the same operators.
Hence the operator layer is single --- one basis, read by all sectors --- and the sector
differences are differences of role, not of operator.
\end{proof}

\begin{corollary}[Organization is a read of the operators]
\label{cor:organization-read}
The organization structure of the theory is a read of the operator basis over the D96
spectrum: the organizational quantities are operator projections.
\end{corollary}

\begin{proof}
By Theorem~\ref{thm:operators-derived}, all named quantities are operator projections. By
Chapter~\ref{chap:d96} (Theorem~\ref{thm:organization-derived}), the organization constants
are derived from the spectrum. Hence organization is the operator read of the spectrum ---
the organizational structure is the operator structure.
\end{proof}

\section{Summary}
\label{sec:operators-summary}

This chapter has established the operator basis of the theory. We have proved:

\begin{enumerate}
\item \textbf{Operators as projections} (Theorem~\ref{thm:operators-derived},
Corollary~\ref{cor:operators-not-primitive}): the operators are derived projections of the
D96 spectrum, not primitives;
\item \textbf{Crowding} (Theorem~\ref{thm:crowding-moments},
Corollary~\ref{cor:crowding-degeneracy}): the degeneracy operator projects the moments;
\item \textbf{Compression} (Theorem~\ref{thm:compression-occupancy},
Corollary~\ref{cor:compression-octave}): the octave-banding operator projects the occupancy
structure;
\item \textbf{Beat} (Theorem~\ref{thm:beat-span}, Corollary~\ref{cor:beat-extent}): the
frequency-ratio operator projects the span;
\item \textbf{Locking} (Theorem~\ref{thm:locking-gap},
Corollary~\ref{cor:locking-separation}): the spectral-gap operator projects the separation
structure;
\item \textbf{Universality} (Theorem~\ref{thm:operators-universal},
Lemma~\ref{lem:operators-nontrivial}): the operators are universal across organized systems
and not trivial statistics;
\item \textbf{Completeness} (Theorem~\ref{thm:four-operators},
Theorem~\ref{thm:no-fifth}, Corollary~\ref{cor:four-search-scoped}): the four-operator
basis is justified, and no fifth operator has been found in any searched domain;
\item \textbf{Consequences} (Lemma~\ref{lem:operators-to-physics},
Theorem~\ref{thm:single-operator-layer}, Corollary~\ref{cor:organization-read}): the
operators bridge spectrum to physics, the operator layer is single, and organization is a
read of the operators.
\end{enumerate}

The operator basis of the theory is therefore the derived reading structure of the D96
spectrum: Crowding (degeneracy), Compression (octaves), Beat (extent), and Locking (gaps),
each an indispensable projection, together universal across organized systems, with no fifth
operator found in any searched domain. The following chapters read the lock structure and
then each physics sector through this operator layer.

\begin{thebibliography}{99}
\bibitem{QG260} AT-QG Phase 260, \emph{Resonance Layer Audit}, Docs/Research.
\bibitem{QG261} AT-QG Phase 261, \emph{Resonance Operator Audit} (Operator Layer), Docs/Research.
\bibitem{QG262} AT-QG Phase 262, \emph{Operator Sector Audit}, Docs/Research.
\bibitem{QG263} AT-QG Phase 263, \emph{Single Resonance Dynamics Audit}, Docs/Research.
\bibitem{QG300} AT-QG Phase 300, \emph{Operator Universality Prediction}, Docs/Research.
\bibitem{QG302} AT-QG Phase 302, \emph{Cross-Domain Universality Audit}, Docs/Research.
\bibitem{QG304} AT-QG Phase 304, \emph{Real Network Universality Audit}, Docs/Research.
\bibitem{QG307} AT-QG Phase 307, \emph{Fifth Operator Search}, Docs/Research.
\bibitem{QG308} AT-QG Phase 308, \emph{Operator Necessity Audit}, Docs/Research.
\bibitem{QG309} AT-QG Phase 309, \emph{Red Team Audit}, Docs/Research.
\bibitem{QG312} AT-QG Phase 312, \emph{Adversarial Spectrum Audit}, Docs/Research.
\bibitem{MONO007} AT-MONO007, \emph{Referee-Readiness Resolution}, Docs/Zenodo/Monograph.
\end{thebibliography}
